%%%%%%%%%%%%%%
% Fichero: uclmTFGesi.tex
% Autor: Jesús Salido Tercero (http://www.uclm.es/profesorado/jsalido)
% Fecha (creación): Febrero 2010 
% Rev. : Febrero 2019
% Descripción: Plantilla para memoria de TFG 
% (Escuela Sup. de Informática, UCLM). Creada para el curso 
% “LaTeX esencial para preparación de TFG, Tesis y otros documentos 
% académicos” (Esc. Sup. Informática-UCLM)
%
% Comentarios: Preparada para `pdflatex' y `biblatex' (con `biber'). 
% Documento editado con TeXstudio. 
% Para su compilación se aconseja utilizar latexmk y biber (requiere perl):
% $latexmk -pdf -silent -synctex=1 --enable-write18 % (última opción innecesaria en Unix)
% En sistemas Unix: latexmk y biber deben ser instalados por separado.
%
% Una versión actualizada de esta plantilla está disponible en overleaf.
% Puede crearse un proyecto propio para escribir un TFG directamente en overleaf,
% o bien descargarla como un archivo .zip para su utilización en modo local.
%%%%%%%%%%%%%%


% -------------------------
%
% PREÁMBULO del documento
% (Editar sólo en caso de cambio del idioma pral. del documento u
%  opciones de los paquetes incluidos).
%
% -------------------------
% EDITAR: Idioma pral. \spanishtrue (Español), \spanishfalse (Inglés)
%Comentar opción no deseada
\newif\ifspanish%Definición de un condicional (spanish)

\spanishtrue 			% Idioma pral.: español (spanish=true)
%\spanishfalse			% Idioma pral.: inglés (spanish=false)
	
\documentclass[ 		% Clase del documento
	11pt,				% Tamaño de letra
	a4paper,			% Tamaño de papel
	twoside,			% Impresión a doble cara
	openright,			% La apertura de cap. a la dcha.
	final       		% Versión final
]{book}

\usepackage[utf8]{inputenx} % Codificación de entrada
\usepackage[english,spanish,es-tabla,es-noindentfirst]{babel} % Internacionalización
%\usepackage{indentfirst} % Para asegurar sangrado en 1ª línea tras sección (necesario con varios idiomas)

%--- Geometría de las páginas del documento
\usepackage[			% Márgenes del documento
	top=2.5cm,			% Margen superior
	bottom=2.5cm,		% Margen inferior
	inner=3.5cm,		% Margen al interior
	outer=2cm			% Margen al exterior
]{geometry}


%--- Tipografía
\usepackage{textcomp,marvosym,pifont} % Generación de símbolos especiales
\usepackage{ccicons} % Iconos de licencia Creative Commons
\usepackage{amsmath,amsthm,amssymb}	% Mejoras cuando hay matemáticas


%--- Tipografía (Opción 1)
\usepackage[tt=false]{libertine}	% Libertine
\usepackage[libertine]{newtxmath}	% Times
%---


%--- Tipografía (Opción 2)
%\usepackage{newpxtext}				% Palatino: La opción osf proporciona números en old style.
%\usepackage{newpxmath}				% Palatino
%---


\usepackage[T1]{fontenc}% Codificación de salida    
\usepackage{microtype}	% Mejoras de microtipografía en la obtención de PDF (sólo para pdflatex)


\usepackage{url} 		% Para escritura de URL
\urlstyle{sf}			% Estilo sans serif para URLs


%--- Definición de colores
% Este paquete debe cargarse antes de ctable.
% Revisar la documentación del paquete para ver los nombres de colores predefinidos.
\usepackage[%
	usenames,
	dvipsnames,
	svgnames,
	x11names,
	table
]{xcolor}


%--- Color especial definido para los hiperenlaces
\definecolor{palered}{rgb}{.8,0,0}


%--- Generación de hiperenlaces
\usepackage[
	pdftex,
	breaklinks,			% Permite que los links ocupen más de una línea
%	hidelinks,			% Oculta el color y borde de los links
% OJO: La opción colorlinks se comenta para evitar un error con el paquete menukeys. 
%	colorlinks=true,	% Pone color en los link o un borde
	linkcolor=palered,	% Color de los links
	anchorcolor=palered, 
	citecolor=palered, 
	filecolor=palered, 
	menucolor=palered, 
	urlcolor=palered,
	bookmarksnumbered=true % Incluye números en bookmarks
]{hyperref}


\usepackage{pdfpages}  % Permite inclusión de páginas de un PDF


%--- Paquetes para listas y organización de texto
\usepackage{paralist}	% Mayor control de listas
\usepackage{multicol}	% Texto en varias columnas


%--- Gráficos y tablas
\usepackage{graphicx}	% Inclusión de figuras
\usepackage{subfigure}	% Inclusión de subfiguras
% EDITAR: Si es necesario cambiar el path para los directorios de figuras.
\graphicspath{{./figs/}}% Path de búsqueda de ficheros gráficos
\DeclareGraphicsExtensions{.pdf,.png,.jpg} % Precedencia de extensiones
\usepackage{rotating}	% Giro de cajas (texto, figuras, tablas) (No DVI)
\usepackage{tabularx,booktabs}	% Ajustes para tablas


%--- Paquetes especiales para Informática
\usepackage{listings}	% Inclusión de listados de código

% Inclusión de algoritmos
\usepackage[
	lined,
	boxruled,
	algochapter,
	commentsnumbered,
\ifspanish
	spanish
\else
	english
\fi
]{algorithm2e}



%--- Personalización de títulos de figuras y tablas
\usepackage[%
	margin=10pt,		% Margen
	font=small,			% Tamaño de tipografía
	labelfont=bf,		% Prefijo-Etiqueta en negrita
	format=hang			%
]{caption}
\captionsetup[table]{skip=4pt} 	% Separación del caption en las tablas


%--- Bibliografía: Biblatex con biber.
% OJO: Para bib multilingüe añadir campos language y langid en registros bib.
\usepackage[
	backend=biber, 		% Backend
	sortcites,
	defernumbers=true, 	% Para numerar al final
	style=numeric-comp, % Estilo numérico condensado
	% Descomentar las opciones siguientes para bibliografía multilingüe
%	autolang=other, 	% Requerido para opción multilingüe
%	language=auto   	% Requerido para opción multilingüe
]{biblatex}


% Línea añadida para eliminar el idioma de la fuente bibliográfica.
\AtEveryBibitem{\clearfield{note} \clearlist{language}}
% OJO: Editar si se cambia el fichero de bibliografía. 
\addbibresource{biblioTFG.bib} 	% Fichero de bibliografía.
\usepackage[autostyle]{csquotes}
%---



%--- Paquete con personalización local para el TFG (ESI-UCLM)
\usepackage{uclmTFGesi}
% -------------------------
% PAQUETES QUE USA: 
% 	fancyhdr,titlesec,sectsty,tikz
% -------------------------
% LISTA DE COMANDOS PROPORCIONADOS:
% OB-Obligatorio
% OP-Opcional
% RE-Recomendado
%
% Comandos para definir variables con los datos del documento.
%
% OB-\portadaTFG				: Pág. de portada (usa variables definidas)
% OB-\portadillaTFG				: Pág. de portada interna
% OB-\tribunalTFG				: Pág. tribunal
% RE-\dedicado{texto}			: Pág. de dedicatoria con texto
% RE-\creditos{texto}{imagen}	: Pág. de créditos con el texto y la imagen
% OB-\abstract{texto}			: Añade abstract (no definido en clase book)
% OP-\tecla{texto}				: Borde de tecla en torno al texto (sustituir por menukeys)		
% OP-\nodivide[penalty]			: Penaliza la división de palabras. 
%								  Máx. (n=10000) sin arg.
% OP-\nowidowandorphan[penalty]: Penaliza las viudas y huérfanas.
%								(sólo si necesario)
% OP-\nodividenotes[penalty]	: Penaliza la división de notas al pié 
%								entre págs.	(sólo si necesario)	
% OP-\savepagecnt				: Crea contador interno con el nº de pág. actual
% OP-\contpagination			: Recupera el valor de pág. previamente
%								 salvado en el cont. interno.
% RE-\cleanhdfirst				: Elimina la cabecera en la primera 
%								página de capítulo.
%
% -------------------------




%--- Paquete para índice temático
\usepackage{makeidx}
\makeindex


%--- Paquete para incluir menús, paths y teclas de modo "elegante"
% OJO: Este paquete presenta algunas incompatibilidades, debe cargarse el último y exige la desactivación de la opción colorlinks y la inclusión en este paquete de la opción hyperrefcolorlinks.
% Este paquete presenta alguna incompatibilidades por lo que cambiarlo de ubicación puede generar algún error (manejar con cuidado).
\usepackage[os=win,hyperrefcolorlinks]{menukeys}

% Estas definiciones permiten cambiar el estilo de los elementos. Si se desean otros estilos o su configuación es preciso recurrir a la documentación del paquete (no lo recomiendo).
\renewmenumacro{\menu}[>]{menus} % default: menus
\renewmenumacro{\directory}[/]{pathswithblackfolder} % default: paths
\renewmenumacro{\keys}[+]{shadowedroundedkeys} % default: roundedkeys

%%%%%%%%%%%%%%%%%%%%%%%%%%%%%%%%%%%%%%%%%%%%%%%%%%%%%
%%%%%%%%%%%%%%%%%%%%%%%%%%%%%%%%%%%%%%%%%%%%%%%%%%%%%
%%%%%%%%%%%%%%%%%%%%%%%%%%%%%%%%%%%%%%%%%%%%%%%%%%%%%





% -------------------------
%
% DATOS DEL DOCUMENTO 
% Definición de variables empleadas en el documento por lo que no son
% traducidos. Cuando algún campo puede tener varias líneas aparecen dos
% campos señalados como <campo>Primera y <campo>Segunda.
% Si no se desea emplear un campo este debe comentarse.
%
% -------------------------
% EDITAR: Datos del documento
\tituloPrimera{Plantilla de TFG para ESI-UCLM}	% 1ª Línea
\tituloSegunda{Curso de {\LaTeX{}} esencial}			% 2ª Línea  % Comentar si no existe
\titulo{Plantilla TFG}							% Título corto
\autor{Jesús Salido Tercero}
\email{jesus.salido@uclm.es}					
\director{<director (nombre apellidos)>}
\codirector{<codirector (nombre apellidos)>}	% Comentar si no existe
\instEdu{UNIVERSIDAD DE CASTILLA-LA MANCHA}
% Fichero con escudo de la institución
% Poner el logo del centro que corresponda
%\escudo{esi} 								
%\escudo{esicolor} 
%\escudo{logo_ESI} 							
\escudo{escudoInf} 	% Nucleo de ferrita						
%\escudo{uclm} 
%\escudo{etsii} 

\centroEdu{ESCUELA SUPERIOR DE INFORMÁTICA}
\deptoEduPrimera{<Primera línea Depto. Director>}% 1ª Línea
\deptoEduSegunda{<Segunda línea Depto. Director>}% 2ª Línea  % Comentar si no existe
\titulacion{GRADO EN INGENIERÍA INFORMÁTICA}
\especialidad{<Tecnología Específica>} 			% Tecnología específica, ...
\tipoDoc{TRABAJO FIN DE GRADO}
% Si las fechas se desean en inglés hay que ponerlas explícitamente.
\fechaDef{<mes, año>} 							% Fecha de defensa
\mesDef{<mes>}        							% Mes de defensa
\yearDef{<año>}        							% Año de defensa
\lugarDef{Ciudad Real}							% Lugar de defensa


% --- Propiedades para el documento PDF
\hypersetup{%
% EDITAR: Valores para el PDF.	
	pdftitle={Plantilla TFG}, 	% Título
	pdfauthor={J. Salido},	% Autor
	pdfsubject={Plantilla LaTeX},	% Tema
	pdftoolbar=true,		% Muestra la toolbar de Acrobat
	pdfmenubar=true			% Muestra la menubar de Acrobat
}
% -------------------------




% -------------------------
%
% CUERPO DEL DOCUMENTO
%
% -------------------------
\begin{document}
	
%--- Ajustes del documento (en páginas iniciales).
\ifspanish
	\selectlanguage{spanish}% Emplea idioma español
\else
	\selectlanguage{english}% Emplea idioma inglés
\fi

\frontmatter
% Cambia la numeración de páginas a números romanos y las secciones no están numeradas aunque si aparecen en el índice de contenidos.
\pagestyle{empty}  % Páginas sin cabecera ni pies
%---



% -------------------------
%
% PORTADAS
%
% -------------------------
\portadaTFG		% Portada pral.

\portadillaTFG	% Portada interior (añade tutores)
%---






% -------------------------
%
% CRÉDITOS
%
% -------------------------
% EDITAR (opcional): Licencia (si se desea modificar).
% Este comando permite una gran flexibilidad y la ventaja de no depender de paquetes externos.
% Esta es una página reservada para señalar información relativa a los derechos de autor y la licencia de distribución y uso del documento. Esta página debería ser aprovechada también para informar de cualquier tipo de cesión de los derechos anteriormente citados. El autor del TFG debe tener presente que el incumplimiento de la legislación vigente en materia de protección de la propiedad intelectual es de su exclusiva responsabilidad independientemente de la cesión de derechos que este haya convenido para su obra ya que no son objeto de cesión aquellos derechos de los que no se es poseedor.

\creditos{Este documento se distribuye con licencia Creative Commons Atribución Compartir Igual 4.0. El texto completo de la licencia puede obtenerse en \url{https://creativecommons.org/licenses/by-sa/4.0/}. 
% El escudo de Informática basado en el nuclo de ferrita que acompaña la distribución de esta plantilla ha sido realizado por P.~Moya, D.~Villa e I. Díez, su inclusión debe respetar los derechos de autor y las licencias a las que se vea sometido. 

La copia y distribución de esta obra está permitida en todo el mundo, sin regalías y por cualquier medio, siempre que esta nota sea preservada. Se concede permiso para copiar y distribuir traducciones de este libro desde el español original a otro idioma, siempre que la traducción sea aprobada por el autor del libro y tanto el aviso de copyright como esta nota de permiso, sean preservados en todas las copias.}{cclicense}
%---




% -------------------------
%
% TRIBUNAL
%
% -------------------------
\tribunalTFG % Página para calificaciones del tribunal
%---





% -------------------------
%
% DEDICATORIA (opcional, 1 pág. máximo)
% Aunque opcional, no se debería perder la oportunidad de poder 
% dedicar el trabajo a alguien MUY especial.
%
% -------------------------
% EDITAR: Dedicatoria (comentar si no se desea incluir).
\dedicado{A mis alumnos ... \\ % A alguien muy especial
(para siempre)} % Como mucho dos líneas (no confundir con los agradecimientos).
%---









%--- Ajustes del documento.
\pagestyle{plain}	% Páginas sólo con numeración inferior al pie

% -------------------------
%
% RESUMEN:
% OJO: Si es preciso cambiar orden manualmente
%
% -------------------------
%--- Resumen en español
\selectlanguage{spanish} % Selección de idioma del resumen.
\cleardoublepage % Se incluye para modificar el contador de página antes de añadir bookmark
\phantomsection  % OJO: Necesario con hyperref
\pdfbookmark[0]{Resumen}{idx_resumen}% idx_resumen.0 % Bookmark en PDF


\begin{abstract}
% EDITAR: Resumen (máx. 1 pág.)
\begin{center}
\emph{(... versión del resumen en español ...)}
\end{center}
El resumen debe ocupar como máximo una página y en dicho espacio proporcionará información crucial sobre el \emph{`qué'} (problemática que trata de resolver el TFG), el \emph{`cómo'} (metodología para llegar a los resultados) y los objetivos alcanzados.
\end{abstract}
%---










%--- Resumen en inglés
% Abstract
\selectlanguage{english} % Selección de idioma del resumen.
\cleardoublepage
\phantomsection % OJO: Necesario con hyperref
\pdfbookmark[0]{Abstract}{idx_abstract}% idx_abstract.0 % Bookmark en PDF

\begin{abstract}
% EDITAR: Abstract (máx. 1 pág.)

\begin{center}
\emph{(... english version of the abstract ...)}
\end{center}
Versión del resumen en inglés. En los trabajos cuyo idioma principal sea el inglés, el orden de \textsf{Resumen} y \textsf{Abstract} se invertirá.	
\end{abstract}
%---














% Selección de idioma para el resto del documento y adaptación de títulos de sección.
% EDITAR: Sólo si es necesario.
% NOTA: Al cambiar de idioma las def. de títulos se reinician.
\ifspanish
	\selectlanguage{spanish}% Para el resto del documento el idioma es español.
	%--- No necesarios al añadir la opción es-tabla para babel.
	%\renewcommand{\tablename}{Tabla} % Se sustituye 'Cuadro' por 'Tabla'
	%\renewcommand{\listtablename}{Índice de tablas}
	%---
	\renewcommand{\lstlistingname}{Listado}
	\renewcommand{\lstlistlistingname}{Índice de listados}
	\SetAlgorithmName{Algoritmo}{Alg}{Índice de algoritmos}
	% Modififica las macros \algorithmcfname y \listalgoritmcfname
	\renewcommand{\appendixname}{Anexo}
	\renewcommand{\bibname}{Bibliografía}
	\renewcommand{\indexname}{Índice temático}
\else
	\selectlanguage{english}
\fi

% -------------------------
%
% AGRADECIMIENTOS (recomendable máx. 1 pág.)
%
% -------------------------
\cleardoublepage
\phantomsection % OJO: Necesario con hyperref
\pdfbookmark[0]{Agradecimientos}{idx_agrad}% idx_agrad.0 % Bookmark en PDF

\chapter*{Agradecimientos} % Opción con * para que no aparezca en TOC ni numerada

%OJO: Editar
 Aunque es un apartado opcional, haremos bueno el refrán \emph{<<Es de bien nacidos, ser agradecidos>>} si empleamos este espacio es un medio para agradecer a todos los que, de un modo u otro, han hecho posible que el TFG <<llegue a buen puerto>>. Esta sección es ideal para agradecer a familiares, directores, profesores, compañeros, amigos, etc. 
 
 Estos agradecimientos pueden ser tan personales como se desee e incluir anécdotas y chascarrillos, pero nunca deberían ocupar más de una página.

\makeatletter		
\begin{flushright}
	\textit{\@autor}
\end{flushright}
\makeatother % Agradecimientos etc.












% -------------------------
%
% ÍNDICES
%
% -------------------------
\pagestyle{fancy} % Estilo de página ajustado por fancyhdr

% EDITAR: Si alguno de los índices no existe, su inclusión se puede comentar.
\cleardoublepage
\phantomsection % OJO: Necesario con hyperref
\pdfbookmark[0]{Índice general}{idx_toc}% idx_toc.0 % Bookmark en PDF

\tableofcontents  % Índice general

% Todos los listados se han incluido en el índice de contenidos. De modo automático también quedan añadidos a los bookmarks del PDF. Si se desean eliminiar del TOC se pueden comentar el comando \addcontensline.

\cleardoublepage
\phantomsection % OJO: Necesario con hyperref
\addcontentsline{toc}{chapter}{\listfigurename} % Añade la lista de figuras al TOC (también a bookmarks en PDF)
%\pdfbookmark[0]{\listfigurename}{idx_lof}% idx_lof.0 % Bookmark en PDF
\listoffigures    % Índice de figuras (opcional)

\cleardoublepage
\phantomsection % OJO: Necesario con hyperref
\addcontentsline{toc}{chapter}{\listtablename} % Añade la lista de tablas al TOC (también a bookmarks en PDF)
%\pdfbookmark[0]{\listtablename}{idx_lot}% idx_lot.0 % Bookmark en PDF
\listoftables % Índice de tablas (opcional)

\cleardoublepage
\phantomsection % OJO: Necesario con hyperref
\addcontentsline{toc}{chapter}{\lstlistlistingname} % Añade la lista de listados al TOC (también a bookmarks en PDF)
%\pdfbookmark[0]{\lstlistlistingname}{idx_lol}% idx_lol.0 % Bookmark en PDF
\lstlistoflistings % Índice de listados creados con listings (opcional)

\cleardoublepage
\phantomsection % OJO: Necesario con hyperref
\addcontentsline{toc}{chapter}{\listalgorithmcfname} % Añade la lista de algoritmos al TOC (también a bookmarks en PDF)
%\pdfbookmark[0]{\listalgorithmcfname}{idx_loa}% idx_loa.0 % Bookmark en PDF
\listofalgorithms % Índice de algoritmos creados con algortihm2e





























%--- MAINMATTER
% Capítulos del documento
% Salva en un contador interno el nº de páginas actual
% Debe ir antes de \mainmatter (antes de que se reinicie el cnt page)
\savepagecnt
\mainmatter
% Justo antes del primer capítulo del libro. Activa la numeración con números arábigos y reinicia el contador de páginas.

% OJO: No cambiar de ubicación. Elimina cabecera y pie en pág. inicial de cap.
\cleanhdfirst

% Reajuste del número de página consecutivo para no reiniciar paginación en Cáp. 1
%\contpagination % Comentado para reiniciar paginación (pag. 1)




% -------------------------
%
% CAPÍTULOS
%
% -------------------------
% Se incluye un fichero para cada capítulo. Se emplea la instrucción \include porque en los libros lo más habitual es que cada cápítulo comience en una nueva página.
\chapter{Introducción}
\label{cap:Introduccion}

Introducción al tema, entorno y justificación de la importancia del trabajo abordado.

Estructura del resto del documento.

% ------------------------------------------------------------------------------
% Ejemplos para la plantilla
% ------------------------------------------------------------------------------
\section{Ejemplos de listas}
\label{sec:ejListas}
\index{ejemplos} % Véase cómo se incluyen entradas en el índice alfabético
A continuación se van a añadir algunos ejemplos que pueden emplearse al redactar la memoria.

\index{ejemplos!listas} % Para el índice
\noindent Ejemplo de lista con \emph{bullet} especial. 
% Ejemplo: Lista con bullets especiales
% ============
\begin{itemize}
	\item[*] peras
	\item manzanas
	\item[\ding{170}] naranjas
\end{itemize}

\noindent Ejemplo de lista compacta (también se puede emplear el entorno para enumeraciones \emph{compactenum})
% Ejemplo: Lista con balas
% ============
\begin{compactitem}
	\item peras
	\item manzanas
	\item naranjas
\end{compactitem}


\noindent Ejemplo de lista en varias columnas.
% Ejemplo: Listas en varias columnas
% ============
\begin{multicols}{2} % El parámetro es el número de columnas de la lista
	\begin{compactenum}
		\item peras
		\item manzanas
		\item naranjas
		\item patatas
		\item calabazas
		\item fresas
	\end{compactenum}
\end{multicols}


\newpage


\section{Ejemplos de tablas}
\label{sec:ejTablas}
\index{ejemplos!tablas}
A continuación se incluyen algunos ejemplos de tablas hechas con \LaTeX{} y paquetes dedicados.

% Ejemplo: Tabla con macro \cline
% ==========
\begin{table}[htb]%
	\centering
	\caption{Ejemplo de uso de la macro \texttt{cline}}
	\label{tab:cline}
	\begin{tabular}[t]{|r|l|}
		\hline
		7C0 & hexadecimal \\[1cm] % Ejemplo de separación fijada entre líneas
		3700 & octal \\ \cline{2-2}
		11111000000 & binario \\
		\hline \hline
		1984 & decimal \\
		\hline
	\end{tabular}
\end{table}


\noindent Ejemplo de tabla en la que se controla el ancho de la celda.

% Ejemplo: Ejemplo de tabla con control de la anchura de celda.
% ==========
\begin{table}[htb]%
	\centering
	\caption{Ejemplo de tabla con especificación de anchura de columna}
	\label{tab:anchura}
	\begin{tabular}{ | l | l | l | p{5cm} |}
		\hline
		Día & Temp Mín (\textdegree C) & Temp Máx (\textdegree C) & Previsión \\ \hline
		Lunes & 11 & 22 & Día claro y muy soleado. Sin embargo, la brisa de la tarde puede hacer que las temperaturas desciendan \\ \hline
		Martes & 9 & 19 & Nuboso con chubascos en muchas regiones. En Cataluña claro con posibilidad de bancos nubosos al norte de la región \\ \hline
		Miércoles & 10 & 21 & La lluvía continuará por la mañana pero las condiciones climáticas mejorarán considerablemente por la tarde\\
		\hline
	\end{tabular}
\end{table}

\clearpage


\section{Ejemplos de figuras}
\label{sec:ejFiguras}
\index{ejemplos!figuras}

En esta sección se añaden ejemplos de muestra para la inclusión de figuras simples y subfiguras.

% Ejemplo: Ejemplo de inclusión de figura
% ============
\begin{figure}[htb]
	\centering
	\includegraphics[width=0.4\linewidth]{escudoInf}
	\caption[Ejemplo de figura]{Figura vectorial del escudo de la ESI}
	\label{fig:ejFigure}
\end{figure}


\noindent Ejemplo de figuras compuestas por subfiguras.

% Ejemplo: Ejemplo de inclusión de subfiguras
% ============
\begin{figure}[htb]
	\centering
	\subfigure[Gráfico vectorial PDF]{
		\includegraphics[width=0.4\linewidth]{escudoInf}
		\label{fig:escudoColor}
	} 
	\subfigure[Gráfico png]{
		\includegraphics[width=0.4\linewidth]{escudoInfBW}
		\label{fig:escudoBW}
	}
	\caption[Ejemplo de subfiguras]{Ejemplo de inclusión de subfiguras en un mismo entorno}
	\label{fig:ejSubfigures}
\end{figure}


\clearpage


\section{Ejemplos de listados}
\label{sec:ejListados}
\index{ejemplos!listados}

Ejemplos más representativos de inclusión de porciones de código fuente.

% Ejemplo: Listado Java
% ============
\begin{lstlisting}[language=Java,float=ht,caption={[Código fuente en Java]Ejemplo de código fuente en lenguaje Java},label=lst:java]
// @author www.javadb.com
public class Main {    
// Este método convierte un String a
// un vector de bytes

public void convertStringToByteArray() {

String stringToConvert = "This String is 15";      
byte[] theByteArray = stringToConvert.getBytes();        
System.out.println(theByteArray.length);        
}

// argumentos de línea de comandos 
public static void main(String[] args) {
new Main().convertStringToByteArray();
}
}
\end{lstlisting}



\noindent Otro ejemplo.

\begin{lstlisting}[style=C,float=ht,caption={Ejemplo de código C},label=lst:codC]
// Este código se ha incluido tal cual está 
// en el fichero \LaTeX{}
#include <stdio.h>
int main(int argc, char* argv[]) {
puts("¡Hola mundo!");
}
\end{lstlisting}


\noindent Ejemplo de entrada por consola.

\begin{lstlisting}[style=consola, numbers=none]
$ gcc -o Hola HolaMundo.c
\end{lstlisting}
\chapter{Objetivo}
\label{cap:Objetivo}

Introduce y motiva la problemática (i.e.\emph{\ ¿cuál es el problema que se plantea y porqué es interesante su resolución?})

Debe concretar y exponer detalladamente el problema a resolver, el entorno de trabajo, la situación y qué se pretende obtener. También puede contemplar las limitaciones y condicionantes a considerar para la resolución del problema (lenguaje de construcción, equipo físico, equipo lógico de base o de apoyo, etc.). Si se considera necesario, esta sección puede titularse \emph{Objetivos del TFG e hipótesis de trabajo}. En este caso, se añadirán las hipótesis de trabajo que el alumno pretende demostrar con su TFG.

Una de las tareas más complicadas al proponer un TFG es plantear su \textsf{Objetivo}. La dificultad deriva de la falta de consenso respecto de lo que se entiende por \emph{objetivo} de un trabajo de esta naturaleza. En primer lugar se debe distinguir entre dos tipos de objetivo:

\begin{enumerate}
	\item La \emph{finalidad específica} del TFG que se plantea para resolver una problemática concreta aplicando los métodos y herramientas adquiridos durante la formación académica. Por ejemplo, \emph{<<Desarrollo de una aplicación software para gestionar reservas hoteleras \emph{on-line}>>}.
	
	\item El \emph{propósito académico} que la realización de un TFG tiene en la formación de un graduado. Por ejemplo, la \emph{adquisición de competencias específicas de la especialización} cursada.
\end{enumerate}

En el ámbito de la memoria del TFG se tiene que definir el primer tipo de objetivo, mientras que el segundo tipo de objetivo es el que se añade al elaborar la propuesta de un TFG presentada ante un comité para su aprobación. Este segundo tipo de objetivo no debe incluirse en el apartado correspondiente de la memoria y en todo caso puede valorarse su satisfacción en la sección de resultados y conclusiones.

Un objetivo bien planteado para el TFG debe estar determinado en términos del \emph{<<producto final>>} esperado que resuelve un problema específico. Es por tanto un sustantivo que debería ser \emph{concreto} y \emph{medible}. El \textsf{Objetivo} planteado puede pertenecer una de las categorías que se indica a continuación:
\begin{itemize}
	\item \emph{Diseño y desarrollo de <<artefactos>>} (habitual en las ingenierías),
	\item \emph{Estudio} que ofrece información novedosa sobre un tema (usual en las ramas de ciencias y humanidades), y
	\item \emph{Validación de una hipótesis} de partida (propio de los trabajos científicos y menos habitual en el caso de los TFG).
\end{itemize}

Estas categorías no son excluyentes, de modo que es posible plantear un trabajo cuyo objetivo sea el diseño y desarrollo de un <<artefacto>> y éste implique un estudio previo o la validación de alguna hipótesis para guiar el proceso. En este caso y cuando el objetivo sea lo suficientemente amplio puede ser conveniente su descomposición en elementos más simples hablando de \emph{subobjetivos}. Por ejemplo, un programa informático puede descomponerse en módulos o requerir un estudio previo para plantear un nuevo algoritmo que será preciso validar. 

La descomposición de un objetivo principal en subobjetivos u objetivos secundarios debería ser natural (no forzada), bien justificada y sólo pertinente en los TFG de gran amplitud.

Junto con la definición del objetivo del TFG se puede especificar los \emph{requisitos} que debe satisfacer la solución aportada. Estos requisitos especifican \emph{características} que debe poseer la solución y \emph{restricciones} que acotan su alcance. En el caso de TFG cuyo objetivo es el desarrollo de un <<artefacto>> los requisitos pueden ser \emph{funcionales} y \emph{no funcionales}.

Al redactar el objetivo de un TFG se debe evitar confundir los medios con el fin. Así es habitual encontrarse con objetivos definidos en términos de las \emph{acciones} (verbos) o \emph{tareas} que será preciso realizar para llegar al verdadero objetivo. Sin embargo, a la hora de planificar el desarrollo del trabajo si es apropiado descomponer todo el trabajo en \emph{hitos} y estos en \emph{tareas} para facilitar dicha \emph{planificación}.

%El motivo de tal confusión deriva de <<mezclar>> la definición de la finalidad específica del TFG con el propósito académico del mismo.

La categoría del objetivo planteado justifica modificaciones en la organización genérica de la memoria del TFG. Así en el caso de estudios y validación de hipótesis el apartado de resultados y conclusiones debería incluir los resultados de experimentación y los comentarios de cómo dichos resultados validan o refutan la hipótesis planteada.


\chapter{Metodología}
\label{cap:Metodologia}

En este capítulo se debe detallar las metodologías empleadas para planificación y desarrollo del trabajo,
así como explicar de modo claro y conciso cómo se han aplicado dichas metodologías.


Obtención de la BBDD.

Ante la nula información almacenada de una forma estructurada se ha tenido que buscar
 diversas maneras de obtener y almacenar la información, estructurarla y tratarla para
 poder sacar conocimiento de ella.

En la página de la federación internacional de esgrima (FIE) se almacenan los resultados
 de las competiciones de los últimos tiempos. De ahí se puede obtener los resultados
 de cada competición tanto como el ranking general, pasando por la fase de poules
 y acabando con los tablones eliminatorios de estos mismos. Viendo toda esta información
 almacenada se decide extraer y almacenar dicha información. Para la extracción se
 utilizarán técnicas de scrapping web mediante la cual se descarga la página y se
 obtiene su contenido para tratarlo. En este caso navegamos por dicha página y una
 vez obtenida la información la guardamos en nuestra BBDD. Puesto que no está del
 todo estructurada tendremos que ir almacenando dicha información en diferentes BBDD
 y después juntarlas.

Descripción de la BBDD.

La primera BBDD que generaremos será aquella en la que guardaremos la información
 básica de los asaltos. Para ello guardaremos el ID de la competición, el número
 de tablón en el que se efectuó el asalto, ID del primer tirador, ID del segundo
 tirador, tocados obtenidos por el segundo primer tirador y tocados obtenidos por
 el segundo tirador.

\begin{table}[htb]%
  \centering
  \caption{Estructura BBDD asaltos inicial}
  \label{tab:anchura}
  \begin{tabular}{ | l | l | l | l | l | l | }
    \hline
    Nombre de Campo & Tipo de campo & Ejemplo \\ \hline
    CompetitionID & Texto & 2019-64 \\ \hline
    Tableu & Entero & 32 \\ \hline
    Competitor1 & Texto & /fencers/Anna-KOROLEVA-40351/ \\ \hline
    Competitor2 & Texto & /fencers/Kira-KESZEI-49034/ \\ \hline
    ResultCompetitor1 & Texto & V/15 \\ \hline
    ResultCompetitor2 & Texto & D/3 \\
    \hline
  \end{tabular}
\end{table}

El siguiente paso que tendremos que dar será obtener la información de los tiradores
 para ello ser visitará la página correspondiente. Un ejemplo de página que almacena
 la información de un tirador sería el siguiente http://fie.org/es/fencers/Mario-PERSU-31870
 como se puede observar el final de la URL es el mismo que el identificador almacenado en la
 anterior tabla. De modo que explorando la anterior BBDD podremos visitar las páginas de cada
 tirador almacenado y de ese modo generar una nueva BBDD con toda la información de cada uno de ellos.
 En dicha BBDD tendremos su identificador, edad, ranking (si lo tuvieran), nacionalidad, mano usada
 y arma.

\begin{table}[htb]%
  \centering
  \caption{Estructura BBDD tiradores}
  \label{tab:anchura}
  \begin{tabular}{ | l | l | l | l | l | l | }
    \hline
    Nombre de Campo & Requerido & Tipo de campo & Ejemplo \\ \hline
    competitorID & Si & Texto & ADRIANA-MILANO-36467 \\ \hline
    Age & Si & Entero & 21 \\ \hline
    FieRanking & No & Entero & 300 \\ \hline
    HandNess & Si & Texto & Right \\ \hline
    Weapon & Si & Texto & Epée \\
    \hline
  \end{tabular}
\end{table}

Una vez obtenidos todos los datos referentes a los tiradores tendremos que cruzar
 las dos tablas mencionadas anteriormente de modo que tengamos toda la información
 en una sola BBDD. Esta última BBDD tendrá la información de la primera, sustituyendo
 las columnas de ID de cada tirador por su registro correspondiente en la anterior tabla.

De modo que la estructura será la siguiente
\begin{table}[htb]%
  \centering
  \caption{Estructura BBDD asaltos final}
  \label{tab:anchura}
  \begin{tabular}{ | l | l | l | l | l | l | }
    \hline
    Nombre de Campo & Tipo de campo & Ejemplo \\ \hline
    ComptetitionID & Texto & 2019-176 \\ \hline
    TABLEU & Entero & 32 \\ \hline
    C1\_ID & Texto & Sergey-KHODOS-13869 \\ \hline
    C1\_RANKING & Entero & 22 \\ \hline
    C1\_NATIONALITY & Texto & RUSSIA \\ \hline
    C1\_HANDNESS & Texto & Right \\ \hline
    C1\_WEAPON & Texto & Epée \\ \hline
    C2\_ID & Texto & Laurin-EGGENSCHWILER-5966 \\ \hline
    C2\_RANKING & Entero & 34 \\ \hline
    C2\_NATIONALITY & Texto & SWITZERLAND \\ \hline
    C2\_HANDNESS & Texto & Right \\ \hline
    C2\_WEAPON & Texto & Epée \\ \hline
    C2\_ID & Texto & Laurin-EGGENSCHWILER-5966 \\ \hline
    C1\_RESULT & Texto & V/15 \\ \hline
    C2\_RESULT & Texto & D/3 \\
    \hline
  \end{tabular}
\end{table}

Modo de empleo.

El modo de empleo de esta BBDD será para el entrenamiento de una red neuronal
 la cual servirá para apoyar al sistema experto. En total hay unos 28.000 registros,
 de los cuales se emplearán entorno al 60 por ciento para entrenar al modelo y un 40
 por ciento para comprobar que el entrenamiento ha sido satisfactorio.

\chapter{Resultados}
\label{cap:Resultados}

En esta sección se describirá la aplicación del método de trabajo presentado en el capítulo \ref{cap:Metodologia} en este caso concreto, mostrando los elementos (modelos, diagramas, especificaciones, etc.) más importantes. Este apartado debe explicar cómo la metodología satisface los objetivos y requisitos planteados.
\chapter{Conclusiones}
\label{cap:Conclusiones}

Breve resumen de lo más destacable del TFG con la solución propuesta y posibles mejoras, ampliaciones o trabajos relacionados que quedan por hacer y que tienen interés para el tema tratado.


% -------------------------














% -------------------------
%
% ANEXOS
%
% -------------------------
\appendix

% Tras este punto los capítulos se numeran con letras.
% Aquí todos los apéndices necesarios
\chapter{Anexo A}
\label{cap:Test Slagel}


pendiente escribir


 % Apéndice A (opcionales)

%---















%--- BACKMATTER
\backmatter
% -------------------------
%
% BIBLIOGRAFÍA
%
% -------------------------
% OJO: Todas las referencias deben estar citadas en el texto)
% EDITAR: Comentar línea siguiente
\nocite{*} % INCLUIDO para ver cómo queda, pero comentar en versión final.

\phantomsection  % OJO: Ojo necesario con hyperref.
\addcontentsline{toc}{chapter}{\bibname} % Añade la bibliografía al Índice de contenidos.
%---
% Opción 1: Bibliografía con todas las fuentes en un apartado.
%---
\printbibliography
%---

%---
% Opción 2: Bibliografía con secciones separadas.
%---
%\printbibheading
%\printbibliography[heading=subbibliography,type=online,title={Fuentes online}]
%\printbibliography[heading=subbibliography,nottype=online,title={Fuentes no online}]
% -------------------------
















% -------------------------
%
% ÍNDICE TEMÁTICO (opcional)
%
% -------------------------
% CONSEJO: Incluir los comandos mientras se escribe cada capítulo ya que hacerlo al final resulta tedioso.
\cleardoublepage
\phantomsection % OJO: Ojo necesario con hyperref.
\addcontentsline{toc}{chapter}{\indexname} % Añade al Índice de contenidos.

%---
% Impresión de Índice Temático (comentar para eliminar)
%---
\printindex  % Facilitado por makeidx (opcional, si no se usa no se imprime)

\end{document}